\begin{abstract}

Procedural content generation is highly relevant and challenging in the video game industry, a topic that has become even more pertinent with the advent of deep neural models. Particularly, adversarial learning is widely employed for generating 2D game levels. This study utilizes a well-known generative adversarial neural model to create levels for the game Super Mario Bros (1983). To further diversify the generated levels, techniques grounded in evolutionary computing algorithms imposing patterns for the generative model are employed (both single-objective and multi-objective), as well as diverse training sets. Evolutionary algorithms are compared in terms of performance, execution time, adherence to training set patterns, and gameplay, the latter evaluated by an agent well-regarded in the literature. The results demonstrate the ability to generate diverse and playable levels, as well as the capacity to manipulate these levels to meet desired properties and difficulties.

\end{abstract}
